\documentclass[a4paper,12pt]{article}
\usepackage[T1]{fontenc}
\usepackage[utf8]{inputenc}
\usepackage[english]{babel}
\usepackage{times}
\usepackage[left=3cm,right=2cm,top=3cm,bottom=2cm]{geometry}
\usepackage{amsmath,amssymb,amsthm,booktabs,array,microtype}
\usepackage[round,authoryear]{natbib}
\usepackage[hidelinks]{hyperref}
\usepackage{url}
\urlstyle{same}
\hypersetup{pdftitle={JayDB2: Proof-Carrying Transactions under Administrator-Defined Laws},
pdfauthor={Gabriel Jota Lizardo},pdfkeywords={databases, proof-carrying code, Bend, TCC}}
\input{../.build/tcc/numbers.tex}
\newcommand{\code}[1]{\texttt{#1}}
\newcommand{\jay}{JayDB2}
\newtheorem{proposition}{Proposition}
\newtheorem{definition}{Definition}
\setlength{\parindent}{1.25cm}
\setlength{\parskip}{3pt}
\setlength{\emergencystretch}{2em}
\widowpenalty=10000
\clubpenalty=10000
\title{JayDB2: Proof-Carrying Transactions under Administrator-Defined Laws}
\author{Gabriel Jota Lizardo\\[4pt]
\normalsize Pontif\'icia Universidade Cat\'olica de Minas Gerais (PUC Minas)\\
\normalsize Belo Horizonte, Brazil}
\date{2026}
\begin{document}
\maketitle
\begin{abstract}
Applications often enforce access restrictions and data invariants through a
combination of database constraints and application code. This work investigates
a different interface: a database accepts a program together with a proof that
its execution obeys administrator-defined propositions. \jay{} implements this
interface with Bend 2 programs, a private transaction proposal, and a publication
boundary governed by eight fixed laws. We describe its formal model, proof
admission, prepared programs, complete-image persistence, and certified concurrent
reads. Checked source proofs establish the pure kernel and an ordinary serial service model under
explicit toolchain and host-refinement premises. They do not establish an
unconditional theorem about the complete executable. Evaluation combines
positive proofs, deliberately invalid certificates, native execution,
fault-injected recovery, and independently checked benchmark responses. In a
retained two-core experiment, two certified readers achieve median paired
throughput ratios of \TinyRatio{} and \HeavyRatio{} over one reader for two
small-state workloads. These results support an experimental interface for
specialized applications requiring explicit transaction policies, while leaving
substantial proof-authoring, portability, and scaling costs. No general
performance advantage over established database systems is claimed.
\end{abstract}
\noindent\textbf{Keywords:} proof-carrying code; databases; dependent types;
transaction safety; formal verification; Bend.

\section{Introduction}
An application can correctly validate a request in one code path and omit the
same validation in another. Multi-tenant access and relationships among several
tables make this problem particularly visible: protecting each individual row
does not automatically establish that a complete business operation preserves
every intended invariant. Database constraints, transactions and row security
already address important parts of this problem. The question in this work is
whether an executable transaction can instead carry explicit evidence that it
respects the administrator's entire installed policy throughout execution.

\jay{} is an experimental database for that question. An administrator fixes
the schema and laws. A client submits a pure Bend program, a credential,
arguments and a proof. The program describes database effects through a restricted
action language. The service authenticates against published state, checks a
certificate for the actual program and inputs, executes on a private proposal,
and publishes only a complete accepted result. A refusal discards provisional
changes. The caller cannot pass an arbitrary native pointer or directly inspect
the database state.

The research question is: \emph{How can a small programmable database connect
proof admission, effectful execution and atomic publication while making its
remaining implementation assumptions explicit?} A second question concerns
cost: which parts of checking and execution can be reused or overlapped without
weakening that policy boundary? The objectives are to specify the interface,
construct and check source proofs for its kernel, exercise the concrete boundary
with adversarial and fault tests, and measure selected implementations of the
same request semantics.

The contribution is a concrete engineering case study rather than a claim to
have invented proof-carrying code or verified databases. Its useful elements are
the administrator-defined propositions over transaction situations, exact
binding of client certificates to the admitted environment, composition into
eight fixed service obligations, and an executable prototype with reproducible
evidence. Optional indexes and receipts remain ordinary application tables
governed by ordinary laws. This design avoids adding a privileged implementation
for every application abstraction.

\section{Foundations and Related Work}
Proof-carrying code lets the producer of a program supply a machine-checkable
argument that the program respects a consumer's safety policy
\citep{necula1996}. \jay{} applies that general idea to transactions and
administrator propositions. Its certificates target the semantics of Bend
actions rather than an independently verified machine-code semantics. The
difference determines the trusted boundary: source proof acceptance still
requires a correct translation and execution path.

Dependent types let specifications refer to program values. Ynot demonstrates
reasoning about higher-order imperative programs using a dependently typed
setting and an explicit axiomatic extension \citep{ynot2008}. \jay{} uses
propositions expressed as Bend types and a finite action interpreter. This work
does not equate the different logical foundations or transfer Ynot's results to
Bend. Fiat refines abstract data type specifications, including relational query
structures, into implementations with a proof trail \citep{fiat2015}. Its
existence rules out any broad novelty claim that verified query structures or
certified indexes originate here. \jay{} focuses on admitting externally
submitted programs against a fixed live configuration.

Established databases are also stronger comparators than an application with
no checks. PostgreSQL supplies row policies restricting normal row access and
modification, along with integrity constraints \citep{postgresrls,postgresconstraints}.
These mechanisms are valuable alternatives. A logical comparison is whether
the policy appears as runtime predicates or as a proposition certified for a
whole submitted transaction. This is not a measured claim about their relative
security, expressiveness in every setting, or performance. PostgreSQL has a much
broader storage, optimization and operational feature set.

Bend 2 supplies the source language and official checking/compilation tools
\citep{bend}. The prototype pins version 2.0.39 and hashes of its sources and
executables. Pinning makes the evaluated object identifiable. It is not a proof
of the checker, compiler, runtime, or hardware. Later upstream releases and
different environments are outside the measured artifact.

\section{Method}
The study combines artifact construction, deductive source reasoning, negative
controls, and controlled operational experiments. Its unit of correctness is
first the pure transaction interpreter, then an explicitly modeled service, and
finally the concrete host under stated refinement assumptions. These levels
are kept distinct when interpreting results.

The fixed contract is the active \code{contract/} package. Its eight laws and
canonical manifests are inputs, not outcomes edited to make the implementation
pass. The implementation uses source-level proof terms checked in both the
official \code{--check-only} and \code{--verdict} modes. Both modes share a
toolchain and should not be described as two independently implemented trusted
proof assistants. Negative controls remove or contradict required premises and
must fail at the intended obligation. Finite runtime cases provide separate
evidence about compilation and the host bridge.

The principal quantitative dataset is the retained main-review experiment \citep{artifact} at
source revision \code{00d973e58217dbb589e4fdd0a48815a623faae9a}. It identifies
the controller, toolchain, fixture, raw bursts and expected responses. The
later merged validation at \code{1598ed4} records a successful complete check
of the accessory workflow. Measurements from earlier revisions remain
historical; they are not relabeled as fresh measurements of later repairs.
The accompanying evidence script recomputes the reported throughput summaries
and paired ratios from the retained burst records.

This is a single-machine systems case study. There is no user study, external
deployment trial or broad benchmark campaign against production databases.
Correctness experiments intentionally include failure behavior. Performance
experiments deliberately isolate selected volatile read workloads rather than
pretend to estimate all transaction costs.

\section{Data and Transaction Model}
The data model contains schemas, named tables, rows, typed values and a user
registry. Configuration contains schema, laws and technical environment; users
are mutable data. Values distinguish missing rows, optional absence and actual
values. Integer payloads use two 32-bit words representing signed two's-complement
64-bit values. Canonical encoding uses explicit tags, big-endian lengths, valid
UTF-8 and ordered record fields. Rejecting ambiguous representations prevents a
proof and an executable from assigning different meanings to the same message.

Let $C$ denote the fixed configuration, $s$ published state, $c$ an authenticated
context, $v$ request arguments and $p$ a pure client program. The program produces
an action, not a state transformer supplied with unrestricted access to $s$.
Actions describe explicit operations and typed continuations. The interpreter
alone supplies database results. A write returns a unit value, so a caller does
not learn prior row contents or existence through the write result.

Administrator propositions receive situations corresponding to stable state,
an executing prefix, or a closed request. The same complete law list applies in
every phase; phase-sensitive conditions belong inside each proposition. A law
can therefore forbid a foreign-table read even if the program later refuses.
Rollback cannot erase a disclosure that already occurred during an unlawful
prefix. Likewise, a collection helper must justify every actual access and
continuation, not merely its final returned list.

\begin{definition}[Transaction certificate, informal notation]
For the actual $C,c,v,p$, a certificate establishes that, for every valid state
authenticating the same context, all actual execution prefixes obey the installed
laws, actual completion obeys them, and every technically possible publication
preserves required validity. Termination is a separate mandatory obligation.
\end{definition}
This is a readable description of \code{Requirements.Certifies}; the Bend
definition remains normative. The quantification over states prevents a client
from obtaining a proof by assuming a convenient private snapshot. The proof
need not reveal that snapshot to the client compiler.

Consider a task application where each principal owns a table and every row
contains its matching owner. Saving a row requires permission for the actual
table and a correctly owned complete row. Reading and then conditionally
replacing it requires permission for the observation as well as the write.
A refusal after two lawful provisional writes preserves the old published
state. A proof for that workflow does not justify changing an access to another
principal's table. These cases make the policies more useful than an empty
program or a policy that denies every operation.

\section{Architecture and the Eight Laws}
The prototype separates a pure Bend kernel from a Rust host, Node-based source
staging and official-tool invocation, and a small native effect adapter. The
service speaks binary frames on standard input/output. It is not itself a
network database daemon. An optional TypeScript client or web adapter handles
that interface without replacing admission.

Reception checks framing and declared sizes. Authentication uses the current
published registry before the request's program can run. Staging freezes public
configuration and dependencies, constructs targets for the actual source,
context and arguments, checks policy and termination, then compiles the checked
program. Execution occurs in fresh native workers using the private snapshot.
Complete output, child success and cleanup precede acceptance of the proposal.
The host then performs the sole publication and sends the response. An
unsuccessful or incomplete worker cannot supply a complete accepted proposal.

\begin{table}[htbp]\centering\small
\caption{Fixed obligations and their principal source proof components.}
\label{tab:laws}
\begin{tabular}{p{0.07\linewidth}p{0.40\linewidth}p{0.41\linewidth}}\toprule
Law & Required behavior & Representative implementation proof \\\midrule
1 & Valid startup and fixed configuration & \code{Law1Service}, \code{HistoryService} \\
2 & Exact bytes, operations and length checks & \code{CodecLaw2}, \code{Law2Service} \\
3 & Current authenticated identity & \code{Law34Service} \\
4 & Correct certificate bound to the program & \code{Law34Service} \\
5 & Faithful ordered execution and observations & \code{Law5Service}, \code{Machine} \\
6 & Unique valid publication and rollback & \code{Law6Service}, \code{Publication} \\
7 & Exact responses, failures and channel behavior & \code{Law7Service}, \code{ChannelService} \\
8 & Progress under external operating conditions & \code{Law8Service} \\\bottomrule
\end{tabular}
\end{table}

\subsection{Meaning of the checked result}
The checked function \code{ContractService.}\allowbreak\code{contract\_of\_service}
constructs exactly the conjunction
\code{LAWS.Contract(a,k)}. Its assumptions are grouped into initial judgment,
observations, source execution, communication and model availability. They
describe primitive results and their correspondence to the source interpreter,
rather than assume each law as an input. Current-state validity follows from
the initial proof and successive certified publications.

\begin{proposition}[Conditional service correctness]
For the ordinary serial service model, the explicit source/tool/communication premises of
\code{ContractService} entail the eight-law contract for the associated
environment and artifact.
\end{proposition}
\begin{proof}[Proof structure]
The checked construction derives the startup obligation, representation
correctness, authentication and admission obligations, faithful execution,
publication, response/channel obligations and progress. It combines these
terms with the conjunction constructors required by the unchanged contract.
History induction links each call to the preceding publication. The executable
interpretation additionally requires the stated correspondence of concrete
host blocks to this model. The proof terms, rather than this prose sketch,
supply the machine-checked source result.
\end{proof}

The progress law uses the original external \code{Processable} predicate.
Defining processable to mean that this implementation happened to finish would
make progress circular. The source model avoids that substitution. It also
does not promise a maximum elapsed time or physical memory consumption.

The trusted components remain significant: the official checker and
translations, generated native behavior, Rust/Node/C integration, process and
pipe behavior, OS, hardware, and storage assumptions. Neither source hashes nor
many passing tests discharge these premises. Thus ``the Bend proofs check''
and ``every real execution unconditionally obeys all laws'' are different
claims. The latter is not a result of this project.

\subsection{Limits and failures}
The contract has exactly three canonical sizes: message, state and response.
They count encoded bytes and include the boundary value. External frame bytes
are separate. A temporary working state exceeding the final state limit does
not alone justify an early state-size rejection. Finalization determines
whether the completed proposal fits. These limits do not bound the memory of
the complete process tree, compiler work or total execution time.

Response causes distinguish authentication, certificate, refusal, logical
limits and real technical failures. A technical failure cannot substitute for
an arbitrary policy rejection. If delivery fails after publication, the state
remains committed and the caller may not know the result. Fatal uncertainty
must not be represented as a guaranteed rollback. This distinction is crucial
for both durable storage and client retry behavior.

\section{Reusable Programs and Optional Features}
\subsection{Prepared programs}
Ordinary admission may check and compile the same source repeatedly. Prepared
programs instead establish a certificate universally over request contexts and
arguments, compile once, and store an instance-local handle. Execution still
authenticates fresh credentials and binds fresh arguments. The theorem is
specialized to the actual call; it does not authorize a captured old identity.
Restart loses handles. Integrity checks bind the sealed source and native
artifacts to the registration.

Prepared controls extend the wire interface. A raw handle message is not the
ordinary encoded source request in the serial law statement. The source
normalization theorem concerns the resolved logical request; physical framing,
registry integrity and receipt delivery remain explicit host obligations.
This qualification prevents treating a useful adapter proof as a silent proof
of a broader raw protocol. It is a structural boundary: prepared messages use
an initial tag outside the canonical value tags, while the unchanged ordinary
request proposition requires the actual received body to encode an ordinary
request. Prepared-control statuses are also additional interface results.
A whole-wire contract covering these extensions therefore needs an explicitly
approved specification extension or a separately identified ordinary-only
artifact; additional proof terms alone cannot make the raw encodings equal.

\begin{proposition}[Prepared-wire exclusion]
For every ordinary request $r$ and byte-list suffix $t$,
$\operatorname{Encoding}(\operatorname{request\_value}(r),7::t)$ implies
falsehood. If reception processes that same raw body, it has neither an ordinary
\code{RequestReceived} witness nor an ordinary authentication-response
justification.
\end{proposition}
\begin{proof}[Proof structure]
The canonical ordinary request is a record. Encoding either rejects it or emits
record tag 6, which cannot equal the initial tag 7. Transporting this exclusion
through the unchanged reception equation rules out \code{RequestReceived};
the authentication-justification definition requires that witness. Four checked
theorems in \code{PreparedPrograms} construct these implications. This is a
source theorem about the original contract, with no assumed encoding conclusion
or altered rule. Native tests separately establish the prepared host behavior.
\end{proof}

Transparent caching of ordinary full-source requests could retain compilation
reuse without adding wire tags. Client-local handles would expand to source and
proof on every execution. This would defer admission to first execution and
change invalidation semantics. Combined source/proof and arguments can also
exceed the unchanged message ceiling even when the former separate messages
each fit. A no-op wrapper cannot generally inherit a policy that distinguishes
actual program/source/arguments. The compatibility tradeoff therefore needs
an explicit design decision rather than a silent normalization.

\subsection{Persistence and concurrent reads}
Optional persistence stores complete committed images and requires the exact
frozen environment on reopening. The host commits the durable image before
acknowledging success. Recovery depends on a protected publication origin and
exact bytes, not on a checksum magically proving validity. Corrupt, missing or
incompatible storage must not silently become an empty database. An uncertain
commit stops dispatch. Cross-build migration and compatibility are outside the
prototype's scope.

Serial FIFO service is the default. Optional read batches overlap only whole
programs with a separate read-only certificate covering their continuations.
Writes and prepared control operations remain barriers. Readers share the
appropriate published snapshot and responses preserve FIFO order. A source
read-only proof supports this refinement, but the Rust scheduler still needs
its concrete queueing and synchronization contract. More readers can also
increase simultaneous snapshot and worker memory.

\subsection{Libraries, receipts and observation}
The optional library supplies values, records, arithmetic, collection helpers,
queries, conditional mutation and table-based indexes. An index is an ordinary
table with explicit maintenance. Its policy must account for information
disclosure and consistency with primary data. No hidden privilege or implicit
database index maintenance follows from importing a helper.

Receipts pair an operation's exact identity and inputs with a retained result
inside the same transaction as the business effect. Permission to read the
result is checked against current policy. Conflicting reuse must fail. An
explicit retention policy determines when deduplication protection ends.
These mechanisms support retries after a lost reply; a bare transport retry
does not. The TypeScript client therefore does not automatically repeat writes.

An optional administrator timing log measures service work through cleanup and
publication. It excludes queue waiting and response delivery. It changes no
client response and provides no quota enforcement. Observer failures must not
suppress a normal transaction response.

\section{Evaluation}
\subsection{Correctness evidence}
The retained integration review reports checked general, contract, library and
receipt roots in both official modes, with eight aggregate negative controls.
The controls cover missing identity, false byte correspondence, omitted policy,
invalid publication, unguarded allocation length, swapped reification,
authentication/preprocessing order and a hidden write in an alleged read-only
program. Their role is causal: a check that only fails because a dependency is
missing would not establish the intended proof obligation.

On its final reviewed source the retained suite passed 21 admission/verifier,
10 native-runtime, 43 Rust, 19 service, nine recovery and five demo tests.
Application checks exercised 25 indexed durable requests and 43 voucher
requests. Native library tests covered 1,980 signed-integer cases, 1,980
collection cases and 250 variable-length content cases. An additional review
compared 98,305 arithmetic/order results over 19,661 operand pairs and 2,025
composite-key cases against independent oracles. These counts describe finite
checks, not a mathematical proof for arbitrary integer inputs.

Fault tests exercise failures around publication, process death, interrupted
delivery and reopening. They check that acknowledged writes survive, that
uncertain commits stop service, and that application receipts avoid repeating
business writes after a lost response. Such tests cover selected schedules
and injected faults. They cannot exhaust all filesystem, OS or hardware failures.

The subsequent full accessory gate completed with exit zero on an unchanged
checkout, as recorded in \code{docs/accessory-validation.md}. This supplies
evidence for the real TypeScript/persistent-task workflow in addition to the
engine review. A fresh review found further client lifecycle and demo delivery
edge cases and a library aggregate omission. Their repairs and targeted
regressions belong to the current review record. This history illustrates why
a previous passing suite cannot be treated as proof of absence of defects.
The current review also records a complete exit-zero run over 436 unchanged
source, test and configuration inputs. Subsequent pure proof additions have
separate dual-mode checks; the runtime helper implementations remain unchanged
from that run.

\subsection{Read-batch experiment}
The retained comparison uses an Intel Core i7-11700K, Linux x86-64 and two
distinct physical cores. Both configurations receive the same total native
thread budget of two. They differ in one versus two certified read workers.
The client and sampler avoid the service cores and their SMT siblings where
available. Each burst contains eight requests submitted in one pipe write.
There is one warmup burst and seven measured paired rounds per workload and
configuration, giving 256 checked responses including warmup.

Both workloads use a small state containing two users and no application
tables. The tiny workload returns after the read; the heavier workload also
computes a runtime-input recurrence with base $10^8$ dependent U32 updates per
request, alternating $10^8$ and $10^8+1$ across paired rounds. Different seeds
within each burst prevent substituting a single cached result.
An independent arithmetic oracle checks every complete response. Neither
persistence nor timing logs are enabled.

Latency begins just before transmitting an already encoded burst and ends at
receipt of complete responses. It includes queueing, authentication, artifact
checks, worker startup, execution, cleanup and pipe delivery. Client encoding,
expected-result computation and preparation are separate. The recorded setup
cost includes service initialization and registration, so it cannot be labeled
as isolated compilation or proof-checking time.

\begin{table}[htbp]\centering\small
\caption{Retained read-batch measurements. Throughput is the median across seven
bursts. p95 pools measured response latencies. Peak RSS is sampled summed
process-tree RSS, not a memory bound.}
\label{tab:results}
\input{../.build/tcc/table.tex}
\end{table}

The median of the seven paired throughput ratios is \TinyRatio{} for tiny
reads and \HeavyRatio{} for heavier reads. This statistic differs from dividing
two separately computed medians. For tiny reads, the largest paired ratio is
about 3.81, indicating substantial runtime variation. Seven rounds do not
establish a stable population confidence interval or a universal speedup.
The experiment supports overlapping independent requests on this hardware.
It does not show that one serial recurrence becomes faster.

The second reader increases observed memory: the maximum sampled process-tree
RSS rises from approximately 7.05 to 13.71 MiB for tiny reads and from 6.51 to
14.21 MiB for heavier reads. Sampling every 10 ms can miss brief peaks, and
summing RSS double-counts shared pages. Large states, mixed writes, sustained
load and disk behavior require separate experiments. One reader therefore
remains a sensible default.

\subsection{GPU feasibility}
An existing, separate CUDA experiment provides evidence for a narrow extension.
Its focused heavy fixture compares identical complete-request results on CPU
and an NVIDIA RTX 3080 Ti. With 32,768 total arithmetic rounds, median complete
latency is 5.756 ms on CPU and 174.945 ms on CUDA. With 4,294,967,296 rounds,
it is 715.124 ms on CPU and 215.706 ms on CUDA, approximately $3.32\times$
faster for CUDA. The recorded profile uses eight CPU workers, a per-server GPU
worker cap of one and a 512 MiB GPU corpus. These are the historical compared
profiles, not equal CPU/GPU hardware budgets. There are only two complete samples per size/backend and one
separate native trial. These are historical branch measurements, not current
mainline results or a measured general break-even point.

The engineering conclusion is conditional: optional CUDA can make sense for
large explicitly parallel pure computations. General database acceleration
does not follow. Integrating it with the current service would require sealed
GPU artifact identity, authentic current requests, failure rollback, compatible
persistence identity and exclusive GPU scheduling alongside read batches.
The compiler, driver and GPU semantic bridge would remain trusted. Importing
the old branch wholesale would add unsupported combinations. The present
mainline therefore retains CPU execution while recording a concrete future
integration path.

\section{Discussion and Threats to Validity}
The strongest supported safety claim concerns checked source semantics and
the explicitly conditioned service model. It is weaker than an end-to-end
machine-code and operating-system theorem. Conversely, stating the premises
does not make the source proof useless: it exposes the exact places where
operational evidence or a future refinement proof is necessary.

A transaction safety certificate also differs from a functional specification
of every helper. For example, lawful arithmetic may still calculate the wrong
business result, and a lawful read may select the wrong authorized row.
Arithmetic helpers have checked word-level and sign-predicate contracts, but
these do not supply a separate general interpretation theorem over mathematical
integers. Multiplication and division are the largest remaining algebraic gaps,
despite extensive finite tests. The public key
decoder now has a general roundtrip theorem for its actual length-derived fuel,
closing the previous gap between the public function and an internal decoder
given conveniently chosen fuel. Proof coverage must
therefore name the target property and actual function. Counting imports or
law declarations alone cannot measure semantic completeness.

The optional \code{ComparisonProofs} module now establishes soundness and
completeness of structural value equality, including malformed list/record
tails. It also derives strict-order properties, equality reflection and
trichotomy for the actual signed comparator through its sign-flipped word rank.
Constant-true equality and constant-false ordering controls fail at the intended
obligations in both official modes. These results strengthen conditional
comparison and sorting specifications while leaving the separate mathematical
integer interpretation and arithmetic proofs open.

The threat model trusts the administrator, protected host files and the
technical environment. Public configuration and code are separated from the
private seed and state. A hostile same-identity operating-system process or a
defective compiler is not neutralized by source hashing. Credential secrecy,
network authentication, provisioning, side channels, availability budgets and
production operations need their own engineering. Policy authors can also
write a permissive or mistaken proposition: proof establishes the installed
statement, not its adequacy for an unstated intention.

External validity is limited by one machine, small states, selected application
policies and narrow workloads. The performance controls compare implementations
of the same JayDB2 read interface, not SQL engines. No throughput claim against
PostgreSQL or SQLite is justified. A useful next study would compare complete
equivalent workflows, including proof construction and admission, against
database constraints and row policies at several state sizes.

Reproducibility requires more than publishing rounded tables. The repository
retains raw requests/measurements, controller and source identities, expected
results, failure logs and exact commands. The article's generator verifies
counts and recomputes its quantitative table. Nevertheless, historical runs
cannot automatically establish present hardware timing or results for a
modified executable.

\section{Conclusion}
\jay{} demonstrates an executable transaction interface in which clients
supply programs and administrator-policy proofs. The pure kernel and modeled
service connect authentication, certified prefixes, finalization and atomic
publication under explicit integration premises. Real application, negative
admission and recovery tests provide evidence for the concrete prototype.
Prepared programs and optional certified readers reduce selected repeated-work
costs while retaining fresh authentication and the publication boundary.

The most plausible use is a specialized system where a small number of reusable
transactions justify deliberate policy proofs, such as ownership-restricted
tasks, consistent application indexes, or receipt-protected business updates.
The costs include proof expertise, a pinned experimental toolchain,
complete-state transfer and limited operational features. The prototype does
not yet establish that these costs beat established database alternatives.

Future work should first close functional theorem gaps and strengthen the
concrete source-to-executable refinement. Broader state-size and mixed-workload
experiments should precede general scaling claims. Optional GPU execution is
reasonable only with a validated integration and a workload large enough to
amortize startup. The eight laws remain fixed criteria for these improvements.

\begin{thebibliography}{99}\small
\bibitem[Necula and Lee(1996)]{necula1996}
NECULA, G. C.; LEE, P. \textbf{Proof-carrying code}. Pittsburgh: Carnegie Mellon
University, 1996. Technical Report CMU-CS-96-165.
\url{https://personal.utdallas.edu/~hamlen/Papers/necula96.pdf}.

\bibitem[Nanevski et al.(2008)]{ynot2008}
NANEVSKI, A.; MORRISETT, G.; SHINNAR, A.; GOVEREAU, P.; BIRKEDAL, L.
Ynot: dependent types for imperative programs. In: \textbf{ICFP}, 2008,
p. 229--240. DOI: \href{https://doi.org/10.1145/1411203.1411237}{10.1145/1411203.1411237}.

\bibitem[Delaware et al.(2015)]{fiat2015}
DELAWARE, B.; PIT-CLAUDEL, C.; GROSS, J.; CHLIPALA, A.
Fiat: deductive synthesis of abstract data types in a proof assistant.
In: \textbf{POPL}, 2015.
\url{https://hdl.handle.net/1721.1/91993.2}.

\bibitem[PostgreSQL Global Development Group(2026a)]{postgresrls}
POSTGRESQL GLOBAL DEVELOPMENT GROUP. \textbf{PostgreSQL 18 documentation:
row security policies}. 2026.
\url{https://www.postgresql.org/docs/18/ddl-rowsecurity.html}.
Accessed: 6 Nov. 2026.

\bibitem[PostgreSQL Global Development Group(2026b)]{postgresconstraints}
POSTGRESQL GLOBAL DEVELOPMENT GROUP. \textbf{PostgreSQL 18 documentation:
constraints}. 2026. \url{https://www.postgresql.org/docs/18/ddl-constraints.html}.
Accessed: 6 Nov. 2026.

\bibitem[Bend contributors(2026)]{bend}
BEND CONTRIBUTORS. \textbf{Bend}. Source distribution; evaluated toolchain
version 2.0.39, pinned by the JayDB2 admission manifest.
\url{https://github.com/bendlang/bend}. Accessed: 6 Nov. 2026.

\bibitem[Lizardo(2026)]{artifact}
LIZARDO, G. J. \textbf{JayDB2: source, specification and experimental evidence}.
2026. \url{https://github.com/j0taaa/jaydb2}.
Retained measurements: revision \code{00d973e}; merged validation: \code{1598ed4}.
GPU historical fixture: branch \code{experiment-gpu}, revision \code{425074e}.

\bibitem[PUC Minas(2025)]{puc}
PONTIF\'ICIA UNIVERSIDADE CAT\'OLICA DE MINAS GERAIS.
Sistema Integrado de Bibliotecas. \textbf{Orienta\c{c}\~oes para elabora\c{c}\~ao
de projetos de pesquisa, trabalhos acad\^emicos, relat\'orios t\'ecnicos e/ou
cient\'ificos e artigos cient\'ificos}. 6. ed. Belo Horizonte: PUC Minas, 2025.
\url{https://www.pucminas.br/biblioteca}. Accessed: 6 Nov. 2026.
\end{thebibliography}
\end{document}
